\section{Gaussian limits and rank-three examples}
\label{sec:gaussian-benchmarks}

This section gives a sharp Gaussian density bound and explicit examples
for specified input classes or output levels.  For related Gaussian
quadratic-form estimates, see \citet{GNSU2019,BNU2020,CK2026}.

\subsection{The sharp Gaussian ceiling, including affine terms}

\begin{lemma}[Extracting a central square]\label{bench:central-square}
For every \(\lambda\geq1\) and \(m\in\mathbb R\), the law of
\(\lambda(Z+m)^2\), where \(Z\) is standard normal, is a convolution of
the law of a central \(\chi_1^2\) variable and a probability measure
on \([0,\infty)\).
\end{lemma}
\begin{proof}
The elementary Gaussian integral gives, for \(t\geq0\),
\[
 \mathbb E e^{-t\lambda(Z+m)^2}
 =(1+2\lambda t)^{-1/2}
       \exp\left(-\frac{\lambda m^2t}{1+2\lambda t}\right).
\]
First consider the centered factor.  For \(\lambda>1\), the finite
nonnegative measure
\[
 \nu_\lambda(dx)=\frac{e^{-x/(2\lambda)}-e^{-x/2}}{2x}
                 \mathbf1_{\{x>0\}}\,dx
\]
has mass \(\tfrac12\log\lambda\).  A compound Poisson sum with this
intensity measure has Laplace transform
\[
 \exp\left(-\int_0^\infty(1-e^{-tx})\nu_\lambda(dx)\right)
 =\left(\frac{1+2t}{1+2\lambda t}\right)^{1/2}.
\]
The identity follows by integrating
\(\int_0^\infty(e^{-ax}-e^{-bx})\,dx/x=\log(b/a)\);
it may also be checked by differentiation in \(t\) and evaluation at
zero.  At \(\lambda=1\), use the constant zero variable.  Multiplication
by \((1+2t)^{-1/2}\) extracts the central square.

The remaining noncentral factor is the Laplace transform of an
independent compound Poisson sum: take Poisson intensity \(m^2/2\)
and jumps distributed as \(\lambda\chi_2^2\).  Indeed its transform is
\[
 \exp\left[\frac{m^2}{2}\left((1+2\lambda t)^{-1}-1\right)\right]
 =\exp\left(-\frac{\lambda m^2t}{1+2\lambda t}\right).
\]
Uniqueness of Laplace transforms proves the stated convolution identity.
All variables in this distributional construction can be chosen
independently.
\end{proof}

\begin{theorem}[Sharp Gaussian density ceiling]\label{bench:gaussian}
Let \(G\sim N(\mu,\Sigma)\), where \(\Sigma\succeq\sigma^2I\) and
\(\sigma>0\).  If \(A\) is real symmetric and \(s_3(A)\geq a>0\),
then, for arbitrary \(b,c\), the quadratic polynomial
\(G^{\mathsf T}AG+b^{\mathsf T}G+c\) has a continuous density vanishing
at infinity and
\begin{equation}\label{bench:gaussian-bound}
 \|p_{G^{\mathsf T}AG+b^{\mathsf T}G+c}\|_\infty
 \leq\frac{1}{2\sqrt2\,\sigma^2a}.
\end{equation}
The constant is sharp.  In particular, \(\Sigma\succeq I/12\) and
\(s_3(A)\geq1\) give the bound \(3\sqrt2\).
\end{theorem}
\begin{proof}
Write \(G=\mu+\Sigma^{1/2}Z\).  The standard Gaussian vector \(Z\)
may be orthogonally rotated to diagonalize
\(B=\Sigma^{1/2}A\Sigma^{1/2}\).  Its coordinates remain independent.
The singular-value product inequality gives
\[
 s_3(B)\geq s_{min}(\Sigma^{1/2})^2s_3(A)
          \geq\sigma^2a.
\]
For completeness, this inequality follows by writing
\(A=\Sigma^{-1/2}B\Sigma^{-1/2}\) and applying
\(s_j(UDV)\leq\|U\|\|V\|s_j(D)\), which itself follows from the
best rank-\((j-1)\) approximation characterization of \(s_j\).

Put \(d=\sigma^2a\), and divide the output by \(d\).
Select three eigenvalues of \(B/d\) with absolute values at least
one.  Completing the square in each corresponding coordinate and
using Lemma~\ref{bench:central-square} shows that the entire output
is distributed as
\[
 \epsilon_1Z_1^2+\epsilon_2Z_2^2+\epsilon_3Z_3^2+W,
 \qquad \epsilon_j\in\{-1,1\},
\]
where the three displayed standard normals are independent of \(W\).
The constant terms, the extracted Poisson terms with their signs, all
other eigenvalue coordinates, and linear Gaussian terms from the kernel
are absorbed into \(W\).  This argument uses independence after a
Gaussian rotation only.

Convolution cannot increase a density supremum.  If all three signs
coincide, the density is a reflection of the \(\chi_3^2\) density,
whose maximum is
\[
 \max_{x>0}\frac{x^{1/2}e^{-x/2}}{\sqrt{2\pi}}
 =\frac1{\sqrt{2\pi e}}<\frac1{2\sqrt2}.
\]
For the mixed signature, reflection reduces to
\(D=Z_1^2+Z_2^2-Z_3^2\).  The first two squares sum to an exponential
variable of rate \(1/2\).  Therefore, for \(t\geq0\),
\[
 p_D(t)=\frac{e^{-t/2}}{2\sqrt{2\pi}}
                \int_0^\infty v^{-1/2}e^{-v}\,dv
        =\frac{e^{-t/2}}{2\sqrt2}.
\]
For \(u>0\), the same convolution gives
\[
 p_D(-u)=\frac{e^{-u/2}}{2\sqrt{2\pi}}
       \int_0^\infty(u+v)^{-1/2}e^{-v}\,dv
       \leq\frac{e^{-u/2}}{2\sqrt2}.
\]
These formulas give a continuous density, with maximum
\(p_D(0)=1/(2\sqrt2)\), vanishing at infinity.  Convolution with a
probability measure preserves these properties.  Rescaling by \(d\)
proves \eqref{bench:gaussian-bound}.  Equality holds for a centered
isotropic three-dimensional Gaussian and
\(A=a\,\diag(1,1,-1)\).
\end{proof}

\subsection{A global ceiling for diffuse Lindeberg arrays}

\begin{theorem}[Uniform convergence of diffuse quadratic densities]
\label{bench:lindeberg}
For each \(m\), let \(R_m\) be an \(n_m\)-by-\(3\) matrix with
orthonormal columns and rows \(r_{mi}^{\mathsf T}\), such that
\[
 \tau_m:=\max_i\|r_{mi}\|^2\longrightarrow0.
\]
Fix \(\Lambda=\diag(\lambda_1,\lambda_2,\lambda_3)\) with
\(|\lambda_j|\geq1\).  Let the physical coordinates \(X_{mi}\) be
independent, with densities bounded by one, finite means \(\mu_{mi}\),
and finite variances \(v_{mi}\).  Suppose
\begin{equation}\label{bench:lindeberg-limits}
 h_m:=\sum_i r_{mi}\mu_{mi}\longrightarrow h,\qquad
 \Sigma_m:=\sum_i v_{mi}r_{mi}r_{mi}^{\mathsf T}
                        \longrightarrow\Sigma,
\end{equation}
and, for every \(\varepsilon>0\),
\begin{equation}\label{bench:lindeberg-condition}
 \sum_i\mathbb E\left[
   \|r_{mi}(X_{mi}-\mu_{mi})\|^2
   \mathbf1_{\{\|r_{mi}(X_{mi}-\mu_{mi})\|>\varepsilon\}}
                  \right]\longrightarrow0.
\end{equation}
Then the forms
\[
 Q_m=X_m^{\mathsf T}R_m\Lambda R_m^{\mathsf T}X_m
\]
have continuous bounded densities for all sufficiently large \(m\).
If \(G\sim N(h,\Sigma)\) and \(Q_\infty=G^{\mathsf T}\Lambda G\),
then
\begin{equation}\label{bench:uniform-limit}
 \|p_{Q_m}-p_{Q_\infty}\|_\infty\longrightarrow0,
 \qquad
 \limsup_m\|p_{Q_m}\|_\infty\leq3\sqrt2.
\end{equation}
The last constant is sharp within this class.
\end{theorem}
\begin{proof}
Every cap-one density with finite variance satisfies
\(\operatorname{Var}(X)\geq1/12\).  One direct proof, writing
\(\mu=\mathbb EX\), is
\[
 \mathbb P(|X-\mu|>t)\geq1-2t\quad(0\leq t\leq1/2),
 \qquad
 \mathbb E(X-\mu)^2
 \geq\int_0^{1/2}2t(1-2t)\,dt=\frac1{12}.
\]
Consequently \(\Sigma_m\succeq I_3/12\) and
\(\Sigma\succeq I_3/12\).

Here are details of the required vector central limit argument.
Fix \(u\in\mathbb R^3\), and put
\(\xi_{mi}=u^{\mathsf T}r_{mi}(X_{mi}-\mu_{mi})\).
The scalar Lindeberg condition follows from
\eqref{bench:lindeberg-condition}, and
\(\sum_i\mathbb E\xi_{mi}^2\to u^{\mathsf T}\Sigma u\).
The expansion of \(e^{ix}\) to order two, splitting its remainder
at \(|x|=\delta\), gives
\[
 \sum_i\left|
  \mathbb Ee^{i\xi_{mi}}-1+\tfrac12\mathbb E\xi_{mi}^2
            \right|\longrightarrow0.
\]
Indeed, the small part is at most a constant times
\(\delta\sum_i\mathbb E\xi_{mi}^2\), and the large part tends to
zero by Lindeberg; then let \(\delta\downarrow0\).
The same splitting shows
\(\max_i\mathbb E\xi_{mi}^2\to0\).
Taking logarithms of the product of the scalar characteristic
functions therefore proves
\[
 R_m^{\mathsf T}X_m\ \Longrightarrow\ N(h,\Sigma).
\]
The continuity theorem for characteristic functions and the continuous
mapping theorem now give pointwise convergence
\(\varphi_{Q_m}(t)\to\varphi_{Q_\infty}(t)\).

For large \(m\), all row leverages are at most \(1/2\), so
Theorem~\ref{four:diffuse} applies to the actual finite product laws.
Since \(|\det\Lambda|\geq1\), it gives the common integrable bound
\[
 |\varphi_{Q_m}(t)|\leq
 g(t):=\min\{1,(\pi e)^{3/2}|t|^{-3/2}\}.
\]
The pointwise limit satisfies the same bound.  Dominated convergence
and Fourier inversion give
\[
 \|p_{Q_m}-p_{Q_\infty}\|_\infty
 \leq\frac1{2\pi}\int_{\mathbb R}
       |\varphi_{Q_m}(t)-\varphi_{Q_\infty}(t)|\,dt
 \longrightarrow0.
\]
Theorem~\ref{bench:gaussian}, with the covariance floor already proved,
gives the final ceiling.

For sharpness, take \(n_m=3m\), with \(m\) rows equal to
\(e_j/\sqrt m\) for each \(j=1,2,3\), and use independent centered
unit uniforms.  The hypotheses hold with \(h=0\),
\(\Sigma=I_3/12\), and \(\Lambda=\diag(1,1,-1)\).
The limiting density has maximum \(3\sqrt2\), so uniform convergence
gives \(\|p_{Q_m}\|_\infty\to3\sqrt2\).
\end{proof}

\begin{remark}[Hypotheses of the limiting assertion]
The nonzero spectrum of every finite matrix
\(R_m\Lambda R_m^{\mathsf T}\) is exactly the spectrum of \(\Lambda\).
The three projected coordinates need not be independent before the
Gaussian limit.  The proof also allows
\(Q_m=Y_m^{\mathsf T}\Lambda Y_m+\beta_m^{\mathsf T}Y_m+c_m\),
where \(Y_m=R_m^{\mathsf T}X_m\), \(\beta_m\to\beta\), and
\(c_m\to c\), because the physical linear terms are separable phases
in Theorem~\ref{four:diffuse}.  Finite variance, convergence of the
covariances and means, and the Lindeberg condition are hypotheses of
this limiting assertion.
Weak convergence alone would not supply
\eqref{bench:uniform-limit}; the integrable product-law Fourier bound
is the step that excludes narrow surviving density peaks.
\end{remark}

\subsection{One heavy coordinate and two Gaussianized banks}

\begin{theorem}[A uniform gap at the cone tip]\label{bench:one-heavy}
Let \(0<a\leq1\), \(d=\sqrt{1-a^2}\), and
\[
 Q_a=-aX^2+2dXB+aB^2+C^2,
\]
where \(B,C\) are independent \(N(0,1/12)\), independent of \(X\).
Over all probability densities \(0\leq f_X\leq1\), the unique
maximizer of \(p_{Q_a}(0)\), up to null sets, is the centered unit
interval density.  Write the maximum as \(C(a)\).  Then
\begin{equation}\label{bench:one-heavy-bound}
 C(a)\leq
 \frac{\sqrt{6\pi}\,\operatorname{erf}(\sqrt{3/2})+3\sqrt2}{2}
 <4.111374<3\sqrt2.
\end{equation}
This is a bound at output zero.
\end{theorem}
\begin{proof}
Conditional on \(X=x\), complete the square:
\[
 Q_a=a(B+dx/a)^2+C^2-x^2/a.
\]
The density of \((\sqrt a(B+dx/a),C)\) is
\[
 \frac6{\pi\sqrt a}
       \exp\{-6[(u/\sqrt a-dx/a)^2+v^2]\}.
\]
Polar coordinates in \((u,v)\) therefore give the conditional
zero-density kernel, for almost every \(x\),
\begin{equation}\label{bench:heavy-kernel}
 K_a(x)=\frac3{\pi\sqrt a}\int_0^{2\pi}
                      e^{-6F_a(\theta)x^2}\,d\theta,
 \qquad
 F_a(\theta)=\frac{(\cos\theta-d)^2}{a^2}
                         +\frac{\sin^2\theta}{a}.
\end{equation}
The conditional density at any level is at most \(6/\sqrt a\).
For each \(x\ne0\) it is continuous in the level near zero, so
dominated convergence against \(f_X\) proves continuity of the
unconditional density there and identifies its value with
\(\int f_X(x)K_a(x)\,dx\).

The function \(F_a\) is strictly positive.  Thus \(K_a\) is even
and strictly decreasing in \(|x|\).  Filling the interval closest to
zero is consequently the unique cap-one maximizer: explicitly, with
\(I=[-1/2,1/2]\),
\[
 \int (\mathbf1_I-f_X)(K_a-K_a(1/2))\,dx\geq0,
\]
and equality forces \(f_X=\mathbf1_I\) almost everywhere.  Put
\[
 L(s)=\int_{-1/2}^{1/2}e^{-6sx^2}\,dx,
 \qquad R(s)=\sqrt{1+s}\,L(s).
\]
Then \(C(a)=\frac3{\pi\sqrt a}\int L(F_a(\theta))\,d\theta\).
We claim that \(R(0)=1\) and \(R\) is strictly decreasing on
\((0,\infty)\).  If
\(M(s)=\int_{-1/2}^{1/2}x^2e^{-6sx^2}\,dx\), integration by parts
gives \(L(s)-12sM(s)=e^{-3s/2}\), while \(e^y>1+y\) gives
\[
 L(s)=e^{-3s/2}\int_{-1/2}^{1/2}e^{6s(1/4-x^2)}\,dx
       >(1+s)e^{-3s/2}\qquad(s>0).
\]
Since \(L'=-6M\), these identities imply
\[
 \frac{d}{ds}\log R(s)
 =\frac1{2(1+s)}-
        \frac{L(s)-e^{-3s/2}}{2sL(s)}<0.
\]

Write \(\alpha=\sqrt{1+a}\), \(\beta=\sqrt{1-a}\).
Direct expansion gives the perfect-square identity
\[
 a^2(1+F_a(\theta))=(\alpha-\beta\cos\theta)^2.
\]
The tangent-half-angle substitution, or an elementary rational
integral, gives
\[
 \int_0^{2\pi}\frac{d\theta}{\alpha-\beta\cos\theta}
       =\frac{2\pi}{\sqrt{\alpha^2-\beta^2}}.
\]
Consequently
\[
 d\nu_a(\theta)=
 \frac{(1+F_a(\theta))^{-1/2}}{\pi\sqrt{2a}}\,d\theta
\]
is a probability measure and
\begin{equation}\label{bench:heavy-average}
 \frac{C(a)}{3\sqrt2}=\int R(F_a)\,d\nu_a.
\end{equation}

We next show \(\nu_a\{F_a\geq1\}\geq1/2\).
At \(a=1\), \(F_a\equiv1\).  For \(0<a<1\), put
\[
 c_a=\frac{\alpha-a\sqrt2}{\beta},\qquad
 \theta_0=\arccos c_a,\qquad
 \kappa=\sqrt{\frac{\alpha+\beta}{\alpha-\beta}}.
\]
Here \(0<c_a<1\): the upper inequality follows from
\(\alpha+\beta>\sqrt2\), and the lower one also follows from
the comparison below.  The event \(F_a<1\) is exactly
\(\cos\theta>c_a\).  The same substitution gives
\[
 \nu_a\{F_a<1\}
 =\frac2\pi\arctan\left(\kappa\tan\frac{\theta_0}{2}\right).
\]
But
\[
 c_a-\frac\beta\alpha
 =\frac{a(2-\sqrt2\,\alpha)}{\alpha\beta}\geq0.
\]
Equivalently,
\(\kappa\tan(\theta_0/2)\leq1\), so the displayed probability
is at most \(1/2\).  As \(R\leq1\) everywhere and
\(R(s)\leq R(1)\) for \(s\geq1\),
\eqref{bench:heavy-average} yields
\[
 C(a)\leq\frac{3\sqrt2}{2}(1+R(1)),\qquad
 R(1)=\sqrt{\frac\pi3}\,
              \operatorname{erf}(\sqrt{3/2})<1.
\]
This proves the exact bound and its strict gap below \(3\sqrt2\).

The decimal upper bound can also be certified without numerical
quadrature.  Expanding the exponential and integrating gives the
alternating series
\[
 L(1)=\sum_{j=0}^{\infty}
       \frac{(-3/2)^j}{j!(2j+1)}.
\]
Its term magnitudes decrease, so the even partial sum through
\(j=12\) is an upper bound.  Since
\(\sqrt2<14142135623731/10^{13}\), exact rational arithmetic gives
\[
 \frac{3\sqrt2}{2}+3L(1)
 <\frac{3}{2}\frac{14142135623731}{10^{13}}
       +3\sum_{j=0}^{12}\frac{(-3/2)^j}{j!(2j+1)}
 <4.111374.
\]
The expression on the left equals the exact ceiling in
\eqref{bench:one-heavy-bound}.
\end{proof}

\begin{remark}[Exact finite product realization]
The matrix on the three displayed statistics is
\[
 M_a=\begin{pmatrix}-a&d&0\\d&a&0\\0&0&1\end{pmatrix},
 \qquad M_a^2=I_3,\qquad \operatorname{spec}(M_a)=\{1,1,-1\}.
\]
Replace \(B,C\) by two independent normalized sums of \(m\)
centered unit uniforms, retaining the independent heavy coordinate
\(X\).  The corresponding matrix in the \(1+2m\) original
coordinates is \(T_mM_aT_m^{\mathsf T}\), where \(T_m\) has one
row \(e_1^{\mathsf T}\), \(m\) rows \(e_2^{\mathsf T}/\sqrt m\),
and \(m\) rows \(e_3^{\mathsf T}/\sqrt m\).  Its columns are
orthonormal, so its nonzero spectrum is exactly \(\{1,1,-1\}\).
For each fixed \(a>0\) and fixed heavy density, its cone-tip density
converges to that of \(Q_a\).

The bank characteristic function is
\(\psi_m(t)=[\sin(t/(2\sqrt m))/(t/(2\sqrt m))]^m\).
For \(|t|\leq\sqrt m\), the elementary bound
\(|\sin u/u|\leq e^{-c u^2}\) on \(|u|\leq1/2\), for some
\(c>0\), gives a common Gaussian majorant.  For
\(\sqrt m\leq|t|\leq4\sqrt m\), compactness gives the bound
\(\rho^m\), with a fixed \(\rho<1\).  For
\(|t|\geq4\sqrt m\), the bound
\((2\sqrt m/|t|)^m\) has integral tending to zero.  The central
limit expansion on compact intervals and these three bounds prove
\(L^1\) convergence of \(\psi_m\) to the Gaussian characteristic
function.  Thus the bank densities converge uniformly, and are
uniformly bounded, by Fourier inversion.  In the conditional polar
formula used for \eqref{bench:heavy-kernel}, replace the product of
Gaussian bank densities by the product of these two bank densities.
Uniform convergence of the products, the finite angular range, and
\(\int f_X=1\) prove the asserted cone-tip limit.
Neither this argument nor \eqref{bench:one-heavy-bound} gives a
global bound over all output levels, or a convergence rate uniform
as \(a\downarrow0\).
\end{remark}

\subsection{An exact dense Householder example}

\begin{proposition}[The Householder cube cone tip]\label{bench:householder}
Let \(X_1,X_2,X_3\) be independent centered unit uniforms, and put
\[
 e=\frac1{\sqrt3}(1,1,1)^{\mathsf T},\qquad
 A=I-2ee^{\mathsf T}
 =\frac13\begin{pmatrix}1&-2&-2\\-2&1&-2\\-2&-2&1\end{pmatrix}.
\]
Then \(\operatorname{spec}(A)=\{1,1,-1\}\), and the density of
\(Q=X^{\mathsf T}AX\) is continuous at zero, with
\begin{equation}\label{bench:householder-tip}
 p_Q(0)=6\sqrt3\,\operatorname{artanh}
              \left(\frac{\sqrt2-1}{\sqrt3}\right)
       =2.5343552034\ldots.
\end{equation}
In particular, this is also the limit of
\(\mathbb P(|Q|\leq\varepsilon)/(2\varepsilon)\).
\end{proposition}
\begin{proof}
The matrix acts by \(-1\) on \(e\) and by \(+1\) on
\(e^\perp\), proving the spectral statement.  Choose an orthonormal
basis \(u,v\) of \(e^\perp\), and let
\(u_\theta=u\cos\theta+v\sin\theta\).
For \(\mathcal C=[-1/2,1/2]^3\), polar coordinates in the positive
two-plane give the following density version:
\[
 p_Q(s)=\frac12\int_{\mathbb R}\int_0^{2\pi}
  \mathbf1_{\{s+z^2>0\}}
  \mathbf1_{\mathcal C}\bigl(\sqrt{s+z^2}\,u_\theta+ze\bigr)
                                                   \,d\theta\,dz.
\]
Indeed the radial measure is \(r\,dr\), and changing from
\(r\) to \(s=r^2-z^2\) contributes the factor \(1/2\).
Whenever the cube indicator is nonzero, \(|z|\leq\sqrt3/2\).
As \(s\to0\), its integrand converges except at \(z=0\) or
where the limiting ray meets the cube boundary.  For each \(\theta\)
there are only finitely many such nonzero radii.  Dominated convergence
therefore proves continuity at zero.

Put \(d_\pm(\theta)=u_\theta\pm e\) and
\(M_\pm(\theta)=\max_i|d_{\pm,i}(\theta)|\).
At zero, each ray remains in the cube up to
\(r=1/(2M_\pm(\theta))\), so
\[
 p_Q(0)=\frac14\int_0^{2\pi}
        \left(\frac1{M_+(\theta)}+\frac1{M_-(\theta)}\right)d\theta
       =\frac12\int_0^{2\pi}\frac{d\theta}{M_+(\theta)}.
\]
The last equality uses
\(M_-(\theta)=M_+(\theta+\pi)\).
With a suitable origin of the angular coordinate, the three
coordinates of \(d_+\) are phase shifts by \(2\pi/3\) of
\[
 a_0+b_0\cos\theta,\qquad
 a_0=1/\sqrt3,\quad b_0=\sqrt{2/3}.
\]
Every negative coordinate has absolute value at most
\(b_0-a_0<a_0\), whereas the largest positive coordinate is at
least \(a_0\), because their sum is \(\sqrt3\).
Thus the largest absolute coordinate is the largest positive one.
Each coordinate is maximal on a sector of width \(2\pi/3\), giving
\[
 p_Q(0)=\frac32\int_{-\pi/3}^{\pi/3}
                    \frac{d\theta}{a_0+b_0\cos\theta}.
\]
Finally \(t=\tan(\theta/2)\) yields
\[
 \int_{-\pi/3}^{\pi/3}\frac{d\theta}{a_0+b_0\cos\theta}
 =\frac4{\sqrt{b_0^2-a_0^2}}
   \operatorname{artanh}\left(
      \frac1{\sqrt3}\sqrt{\frac{b_0-a_0}{b_0+a_0}}\right).
\]
Since \(\sqrt{b_0^2-a_0^2}=1/\sqrt3\) and
\(\sqrt{(b_0-a_0)/(b_0+a_0)}=\sqrt2-1\), this proves
\eqref{bench:householder-tip}.  Continuity proves the stated
small-ball limit.
\end{proof}

The Householder calculation concerns the cone tip and does not locate
the global maximum of that example.  The one-heavy calculation has the
same output-level restriction.  The diffuse Lindeberg theorem is global
in the output level, under its stated array hypotheses.  These distinctions
leave the unrestricted indefinite product-law endpoint open.
